\section{Related Work}
\paragraph{Classical-to-quantum development.}
C2$|$Q$\rangle$ provides a modular path from classical specifications to encoding,
quantum execution, and decoding~\cite{c2q}. It is an especially relevant comparison
because it begins with classical problem descriptions. Our intended emphasis is
on candidate regions inside surrounding software, the distinction between a
conditional mapping and adoption, and source-linked behavioral obligations. Some
development cores derive from its dataset, so the source dependence is explicit.
This paper does not assert that prior work lacks every form of semantic checking.
Any future direct system comparison must align task scope, access, and contracts.

\paragraph{Quantum code-generation benchmarks.}
Qiskit HumanEval evaluates generated Qiskit programs on curated tasks~\cite{qhe}.
QuanBench combines functional checking with quantum semantic evaluation~\cite{quanbench}.
QuanBench+ extends evaluation across frameworks and includes feedback-based
repair~\cite{quanplus}. The proposed distinction is source-contract evidence for
the same model and separate measurement of claim correctness and completion,
rather than the addition of an execution-feedback loop alone.

\paragraph{Quantum algorithms and workload evaluation.}
Grover search, QAOA, and Ising formulations provide the algorithmic and mathematical
basis of the supported families~\cite{grover,qaoa,lucas}. Their premises must be
connected to source behavior rather than assumed from a problem label.
SupermarQ emphasizes application-oriented quantum benchmarking~\cite{supermarq},
while MQT Bench supports multiple circuit abstraction levels~\cite{mqt}.
These motivate explicit evaluation boundaries and representation-aware resource
reporting. Neither a device benchmark score nor a circuit count is substituted
for preservation of a classical application's observations.

\paragraph{Feedback and synthesis.}
Self-Refine studies iterative improvement using model-generated
feedback~\cite{selfrefine}. Counterexample-guided synthesis separates a proposed
candidate from validation and uses failures to guide subsequent
candidates~\cite{sketching}. The planned workflow adopts a related separation
between model proposal and external evidence, without claiming a general synthesis
algorithm or inheriting formal guarantees. Its research question is whether
program-specific facts and scoped counterexamples improve quantum opportunity
analysis beyond matched self-review. Establishing that contribution requires
the planned controlled study; the current local prototype alone is insufficient.
